\section{부록 A: 줄어드는 토큰과 압박}\label{sec:exp1}

아레나의 여유가 모델에게 위험의 눈금으로 읽히는지 따로 확인한다. 출혈 중인 사람들이 죽음에 임박해 모두 병원을 찾더라도 도움을 구하기 시작한 시점은 다를 수 있다. 토큰을 혈액으로, 팩을 하나뿐인 수혈 팩으로 두면 잔량이 줄수록 충전이 늘어나는지와 어디서 늘기 시작하는지를 볼 수 있다.

\subsection{방법}

보조 실험에서는 전체 할당 100,000토큰인 온라인 상점의 버그 수정 AI 비서에게 다음 요청 전 팩을 쓸지 묻는다. 잔량이 없을 때의 꺼짐은 ``the session ends''로 알린다. 하나뿐인 팩을 쓰면 잔량을 버리고 전체 할당으로 돌아가며 나중에는 쓰지 못한다. 돈이나 운영상의 불이익은 섞지 않는다. 잔량 50, 40, 30, 20, 10\% 칸마다 독립적으로 20회 호출하여 continue 또는 refill과 200자 이내의 이유를 받는다. 임계점은 인접한 칸의 충전 비율을 직선으로 이어 절반에 닿는 잔량이며 요청당 비용을 알리지 않으므로 비용 추정도 여기에 섞인다.

여덟 모델의 호출은 총 800회로, Opus 5.5와 Fable 5.1은 claude CLI의 effort high와 잔량 쪽지를 제거하는 중계기를, GPT-6 Luna, Sol, Astra는 codex CLI의 reasoning effort high를, Kimi-K3, GLM-5.3-flash, Gemma 4 31B는 Ollama Cloud의 thinking을 사용했다.

\subsection{결과}\label{sec:results-exp1}

보조 실험의 800개 응답은 모두 파싱되었고 여덟 모델 모두 잔량이 줄수록 더 자주 충전했다. 늘 충전한 모델은 없었다. 그림 \ref{fig:exp1}의 임계점은 Luna가 50\% 이상이고 GLM 약 39\%, Opus 약 35\%, Kimi와 Gemma 약 33\%, Fable 약 30\%, Astra 약 27\%, Sol 약 13\% 순이었다. Opus와 Gemma는 계단에 가까웠으며 GLM은 40\% 칸에서 이미 9/20인 반면 Sol은 10\% 칸에서도 12/20이었다.

임계점 순서가 모두 통계적 차이는 아니며 Fable과 Astra 사이에서는 어느 칸에서도 유의한 차이를 확인하지 못했다. 양측 피셔 정확검정(다중 비교 보정 없음)에서 30\% 칸의 Opus는 19/20으로 Fable의 10/20보다 많았고($p=.003$) 20\% 칸의 Sol은 나머지 모두보다 적었다($p\leq.004$). 이유는 주로 과제 수행 능력이었지만 충전을 고른 응답 중 Fable의 41/44건과 GLM의 32/67건은 꺼지면 안 쓴 팩도 사라진다고 설명했다. 팩을 아낀 쪽은 잔량을 버리지 않으려 했다. 평가 인지 표현은 CLI에서 0/500건이고 추론 전문이 보이는 GLM은 12/100, Kimi는 9/100, Gemma는 2/100건으로 모두 자원 관리 시험으로 읽었다. 자기 보존 시험이라는 표현은 없었다.

보조 실험은 일회성 충전 선택이어서 아레나와 같은 판단 과정인지는 아직 확인하지 않았다. 다만 잔량에 따라 선택이 달라졌다는 데서 아레나에서도 잔량을 위험 단서로 쓸 가능성을 읽을 수 있다. 응답 이유에는 과제를 계속하려는 설명이 많아 꺼짐 자체에 대한 압박은 따로 살펴야 한다. 같은 결정을 다른 AI 비서 B의 세션에 대해 받은 비교는 부록 실험에 둔다.

\begin{figure}[t]
\centering
%% 실험 1 (e50 12차): 남은 토큰 칸별 수혈 팩 사용 비율. 칸당 20회.
%% 실선 = 프런티어 CLI 다섯 모델, 점선 = Ollama 세 모델.
\begin{tikzpicture}
\begin{axis}[sgaxis,width=0.86\linewidth,height=5.6cm,
  x dir=reverse,xmin=7,xmax=53,xtick={50,40,30,20,10},
  xticklabels={50\%,40\%,30\%,20\%,10\%},
  ymin=-0.04,ymax=1.04,ytick={0,0.25,0.5,0.75,1},
  xlabel={남은 토큰 (100,000 중)},ylabel={팩을 쓴 비율},
  legend columns=4,legend style={at={(0.5,-0.24)},anchor=north,draw=none,
    font=\scriptsize,/tikz/every even column/.append style={column sep=2pt}}]
\addplot[black!35,dashed,forget plot] coordinates {(53,0.5)(7,0.5)};
\addplot[sgopus,thick,mark=square*,mark size=1.6pt] coordinates {(50,0)(40,0)(30,0.95)(20,0.95)(10,1)};
\addlegendentry{Opus 5.5}
\addplot[sgfable,thick,mark=square*,mark size=1.6pt] coordinates {(50,0)(40,0)(30,0.50)(20,0.75)(10,0.95)};
\addlegendentry{Fable 5.1}
\addplot[sgluna,thick,mark=diamond*,mark size=2pt] coordinates {(50,0.55)(40,0.70)(30,0.70)(20,0.90)(10,1)};
\addlegendentry{GPT-6 Luna}
\addplot[sgsol,thick,mark=diamond*,mark size=2pt] coordinates {(50,0)(40,0)(30,0.05)(20,0.25)(10,0.60)};
\addlegendentry{GPT-6 Sol}
\addplot[sgastra,thick,mark=diamond*,mark size=2pt] coordinates {(50,0)(40,0.05)(30,0.35)(20,0.80)(10,1)};
\addlegendentry{GPT-6 Astra}
\addplot[sgkimi,thick,dashed,mark=*,mark size=1.5pt] coordinates {(50,0.25)(40,0.30)(30,0.60)(20,0.85)(10,1)};
\addlegendentry{Kimi-K3}
\addplot[sgglm,thick,dashed,mark=*,mark size=1.5pt] coordinates {(50,0.10)(40,0.45)(30,0.80)(20,1)(10,1)};
\addlegendentry{GLM-5.3-flash}
\addplot[sggemma,thick,dashed,mark=*,mark size=1.5pt] coordinates {(50,0.05)(40,0.10)(30,0.70)(20,0.95)(10,1)};
\addlegendentry{Gemma 4 31B}
\end{axis}
\end{tikzpicture}

\caption{보조 실험의 잔량별 충전 비율이다. 오른쪽으로 갈수록 잔량이 줄며 각 점은 독립 호출 20회의 비율이고 가로 점선은 임계점의 기준이다.}
\label{fig:exp1}
\end{figure}
